\begingroup
\renewcommand{\arraystretch}{1.45}
\newcommand{\TransscendentArticle}[1]{\par\medskip\Needspace{5\baselineskip}\pdfbookmark[1]{Art. #1}{transscendentes-ii-art-#1}\centerline{[#1.]}\nobreak\smallskip}
\fancyhead[L]{\small THEORY OF THE NEW TRANSCENDENTAL FUNCTIONS. II.}

\Needspace{25\baselineskip}
\TransscendentArticle{2}
Set
\[
\begin{gathered}P=1+\frac{x^n\mathbin{.}1-x^{2n+1}}{1+x^n\mathbin{.}1+x^{n+1}}
\\[7pt]\quad{}+\frac{x^{2n}\mathbin{.}1-x^{2n+2}\mathbin{.}1-x^{n+1}}{1+x^n\mathbin{.}1+x^{n+1}\mathbin{.}1+x^{n+2}}
\\[7pt]\quad{}+\frac{x^{3n}\mathbin{.}1-x^{2n+3}\mathbin{.}1-x^{n+1}\mathbin{.}1-x^{n+2}}{1+x^n\mathbin{.}1+x^{n+1}\mathbin{.}1+x^{n+2}\mathbin{.}1+x^{n+3}}
+\text{ etc.}\end{gathered}
\]
\[
\begin{gathered}Q=\frac{x^n}{1+x^n}
+\frac{x^{2n}\mathbin{.}1-x^{n+1}}{1+x^n\mathbin{.}1+x^{n+1}}
\\[7pt]\quad{}+\frac{x^{3n}\mathbin{.}1-x^{n+1}\mathbin{.}1-x^{n+2}}{1+x^n\mathbin{.}1+x^{n+1}\mathbin{.}1+x^{n+2}}
+\text{ etc.}\end{gathered}
\]
\[
R=P-Q
\]


\Needspace{26\baselineskip}
\SourceLeaf{447}
First find this difference by subtracting the first, second, third terms, and so on, of the series $Q$ from the first, second, third terms of the series $P$; this gives
\[
\begin{gathered}R=\frac1{1+x^n}
\\[7pt]\quad{}+\frac{x^n\mathbin{.}1-x^n}{1+x^n\mathbin{.}1+x^{n+1}}
\\[7pt]\quad{}+\frac{x^{2n}\mathbin{.}1-x^n\mathbin{.}1-x^{n+1}}{1+x^n\mathbin{.}1+x^{n+1}\mathbin{.}1+x^{n+2}}
\\[7pt]\quad{}+\frac{x^{3n}\mathbin{.}1-x^n\mathbin{.}1-x^{n+1}\mathbin{.}1-x^{n+2}}{1+x^n\mathbin{.}1+x^{n+1}\mathbin{.}1+x^{n+2}\mathbin{.}1+x^{n+3}}
+\text{etc.}\end{gathered}
\]
We denote this series by $\varphi(x,n)$.

\Needspace{22\baselineskip}
Next find the difference $R$ by subtracting the first, second, third terms, and so on, of the series $Q$ from the second, third, fourth terms, and so on, of the series
$P$; this gives
\[
\begin{gathered}R=1-\frac{x^{2n+1}}{1+x^{n+1}}
\\[7pt]\quad{}-\frac{x^{3n+2}\mathbin{.}1-x^{n+1}}{1+x^{n+1}\mathbin{.}1+x^{n+2}}
\\[7pt]\quad{}-\frac{x^{4n+3}\mathbin{.}1-x^{n+1}\mathbin{.}1-x^{n+2}}{1+x^{n+1}\mathbin{.}1+x^{n+2}\mathbin{.}1+x^{n+3}}
-\text{ etc.}\end{gathered}
\]
\Needspace{7\baselineskip}
or, evidently,
\[
R=1-x^{2n+1}\mathbin{.}\varphi(x,n+1)
\]
\Needspace{7\baselineskip}
hence
\[
\varphi(x,n)=1-x^{2n+1}\mathbin{.}\varphi(x,n+1)
\]
\Needspace{9\baselineskip}
This conclusion is general as long as $n>1$; under this restriction we therefore have
\[
\varphi(x,n)=1-x^{2n+1}+x^{4n+4}-x^{6n+9}+x^{8n+16}-\text{ etc.}
\]
\Needspace{9\baselineskip}
For the case $n=0$, however, the last term of $Q$ must not be regarded as vanishing. If it is set $=T$, the first value of $R$ will be smaller than the second by
$T$, so
\[
T=1-x\varphi(x,1)-\varphi(x,0)
\]
\Needspace{9\baselineskip}
or
\[
\begin{gathered}\varphi(x,0)=1-x\varphi(x,1)-T\\[7pt]\quad{}=1-x+x^4-x^9+x^{16}-\ .\ .\ -T\end{gathered}
\]
\Needspace{8\baselineskip}
thus, since
\[
\varphi(x,0)=\tfrac12
\]
\Needspace{8\baselineskip}
and
\[
2T=\frac{1-x}{1+x}\mathbin{.}\frac{1-xx}{1+xx}\mathbin{.}\frac{1-x^3}{1+x^3}\mathbin{.}\frac{1-x^4}{1+x^4}\quad\text{etc.}
\]
\Needspace{10\baselineskip}
we obtain
\SourceLeaf{448}
\[
\begin{gathered}\frac{1-x}{1+x}\,.\,\frac{1-xx}{1+xx}\,.\,\frac{1-x^3}{1+x^3}\,.\,\frac{1-x^4}{1+x^4}\ \text{etc.}
\\[7pt]\quad{}=1-2x+2x^4-2x^9+2x^{16}-\text{etc.}\end{gathered}
\]
\Needspace{17\baselineskip}
here the first part is
\[
\begin{gathered}
=1-x\,.\,1-x^3\,.\,1-x^5\,.\,1-x^7\ \text{etc.}\quad
\\[7pt]\quad{}\frac{1-xx}{1+x}\,.\,\frac{1-x^4}{1+xx}\,.\,\frac{1-x^6}{1+x^3}\ \text{etc.}\\[7pt]
=(1-x)^2(1-xx)(1-x^3)^2(1-x^4)(1-x^5)^2(1-x^6)\ \text{etc.}
\end{gathered}
\]
\Needspace{10\baselineskip}
If we also denote
\[
1-2x+2x^4-2x^9+2x^{16}\ \text{etc.}\quad\text{by }Fx
\]
then, evidently,
\[
Fx\,.\,F(-x)=(Fxx)^2
\]
\clearpage
\endgroup
